% Đã cấu hình để biên dịch trực tiếp bằng pdfLaTeX (Mặc định của Overleaf)
\documentclass[aspectratio=169]{beamer}

% Cài đặt tiếng Việt cho pdfLaTeX
\usepackage[utf8]{inputenc}
\usepackage[T5]{fontenc}

\usetheme{Madrid}
\usecolortheme{default}

% Packages cho biểu đồ và code
\usepackage{graphicx}
\usepackage{listings}
\usepackage{xcolor}
\usepackage{booktabs}

% Cấu hình hiển thị code JSON
\lstset{
    basicstyle=\ttfamily\scriptsize,
    keywordstyle=\color{blue},
    stringstyle=\color{green!50!black},
    commentstyle=\color{gray},
    showstringspaces=false,
    breaklines=true
}

% Thông tin trang bìa
\title[GLPO AI Optimizer]{GLPO AI Optimizer: Kiến trúc \& Thuật toán Cốt lõi}
\subtitle{Hệ thống Tối ưu hóa Tiến độ \& Chi phí Dự án bằng Trí tuệ Nhân tạo}
\author{Báo cáo Đồ án}
\institute[Đại học]{Khoa Công nghệ Thông tin}
\date{\today}

\begin{document}

% Slide Bìa
\begin{frame}
    \titlepage
\end{frame}

% -------------------------------------------------------------
% PHẦN 1: BÀI TOÁN & THIẾT KẾ DATABASE
% -------------------------------------------------------------
\section{Bài toán \& Kiến trúc}

\begin{frame}{1. Bài toán Thực tiễn: "Sự biến dị của Dữ liệu" (Schema Drift)}
    \textbf{Thực trạng Dataset DSLIB (2011 - 2025):}
    \begin{itemize}
        \item Cấu trúc file Excel dự án thay đổi liên tục qua từng năm.
        \item Các cột dữ liệu co giãn (từ 17 cột xuống 14 cột), xuất hiện nhiều thuộc tính (features) mới.
    \end{itemize}
    \vspace{0.3cm}
    \textbf{Giải pháp Kiến trúc Lai (Hybrid SQL-NoSQL):}
    \begin{itemize}
        \item \textbf{Cột tĩnh (Hard columns):} Lưu trữ thông số cốt lõi (ID, Duration, Fixed Cost).
        \item \textbf{Cột động (JSONB metadata):} Tự động "hút" mọi thông số dị biệt theo từng năm.
        \item Hệ thống phân rã thành \textbf{3 Trục Ràng Buộc} độc lập: Thời gian (Agenda), Tài nguyên (Resource), và Logic (Topology).
    \end{itemize}
\end{frame}

% -------------------------------------------------------------
% PHẦN 2: MIỀN GIÁ TRỊ LÝ THUYẾT & ERD (THEORETICAL DOMAINS)
% -------------------------------------------------------------
\section{Miền Giá trị \& ERD}

\begin{frame}{2. Sơ đồ Thực thể Liên kết (ERD) - Universal Schema}
    Kiến trúc Database được đập đi xây lại hoàn toàn để hỗ trợ Dataset 2011-2025:
    \vspace{0.2cm}
    \begin{columns}
        \column{0.5\textwidth}
        \textbf{Nhóm Bảng Cốt lõi (Core \& Snowflake):}
        \begin{itemize}
            \item \texttt{app\_projects}: Chứa Metadata JSONB.
            \item \texttt{fact\_tasks}: Hút mọi thuộc tính dị biệt.
            \item \texttt{dim\_topic\_cost}: Tách bạch Fixed/Resource.
        \end{itemize}
        
        \vspace{0.2cm}
        \textbf{Nhóm Bảng Ràng Buộc (Constraints):}
        \begin{itemize}
            \item \texttt{project\_constraint\_time} (Lịch trình)
            \item \texttt{project\_constraint\_resource} (Tài nguyên)
            \item \texttt{project\_constraint\_logic} (Mạng DAG)
        \end{itemize}
        
        \column{0.5\textwidth}
        \centering
        \includegraphics[width=0.95\linewidth]{erd.png}
    \end{columns}
\end{frame}

\begin{frame}{3. Miền Giá trị Lý thuyết (Theoretical Domains)}
    Việc định nghĩa Miền xác định giúp AI không đưa ra đề xuất "phi thực tế", là cơ sở cho Constraint Validation.
    \vspace{0.2cm}
    \begin{table}
        \centering
        \footnotesize
        \begin{tabular}{@{}llp{8cm}@{}}
            \toprule
            \textbf{Feature} & \textbf{Domain} & \textbf{Giải thích Toán học} \\
            \midrule
            \textbf{Duration} & $(0, +\infty)$ & Thời lượng không thể âm hoặc bằng 0. \\
            \textbf{Total Float} & $(-\infty, +\infty)$ & Khoảng dự trữ. Có thể âm nếu dự án đã trễ hạn chót (Hard Deadline). \\
            \textbf{Max Availability} & $\mathbb{N} (0, +\infty)$ & Sức chứa tài nguyên là số tự nhiên (Không thể có âm 2 người thợ). \\
            \textbf{Cruciality (CRI)} & $[-1, 1]$ & Hệ số tương quan Pearson đo độ nhạy của Task với toàn dự án. \\
            \bottomrule
        \end{tabular}
    \end{table}
\end{frame}

\begin{frame}{3. Phân rã Đặc trưng theo Cấp độ (Features Range)}
    \textbf{Tầng 1: Task-Level Features (Đặc trưng Cục bộ)}
    \begin{itemize}
        \item \textit{Duration, Resource Cost, Fixed Cost, Predecessors, Optimistic/Pessimistic.}
        \item Đặc điểm: Sẽ biến thiên khi AI áp dụng Crashing (Ép tiến độ).
    \end{itemize}
    \vspace{0.2cm}
    \textbf{Tầng 2: Project-Level Constraints (Ràng buộc Vĩ mô)}
    \begin{itemize}
        \item \textit{Logic Constraint:} Mũi tên phụ thuộc (FS, SS), Lag Days.
        \item \textit{Resource Constraint:} Sức chứa (Max Availability), Loại tài nguyên (Renewable/Consumable).
        \item \textit{Agenda Constraint:} Cấu trúc lịch làm việc \texttt{weekly\_schedule} (VD: Nghỉ Thứ 7, CN).
    \end{itemize}
\end{frame}

% -------------------------------------------------------------
% PHẦN 3: FEATURE ENGINEERING
% -------------------------------------------------------------
\section{Kỹ nghệ Đặc trưng}

\begin{frame}{4. Kỹ nghệ Đặc trưng (Feature Engineering cho AI)}
    Biến đổi dữ liệu thô (Raw Data) thành Vector chuẩn hóa (Scale-Invariant) để Graph Neural Network học:
    \vspace{0.3cm}
    \begin{itemize}
        \item \textbf{Trục Thời gian (Criticality):} \texttt{float\_ratio} = Total Float / Duration. Tỷ lệ tiến về 0 $\rightarrow$ Đường găng (Tử huyệt). Cờ \texttt{is\_critical}.
        \item \textbf{Trục Tài nguyên (Pressure):} \texttt{utilization\_ratio} = Daily Usage / Max Availability. Chuẩn hóa dải $[0, 1]$.
        \item \textbf{Trục Chi phí (Sensitivity):} \texttt{cost\_density} = Resource Cost / Duration. AI dựa vào đây để săn các Task "rẻ nhất" để nén tiến độ.
        \item \textbf{Trục Topological (Hub):} \texttt{descendant\_count}. Đếm số lượng công việc bị kéo lùi nếu Task hiện tại trễ.
    \end{itemize}
\end{frame}

% -------------------------------------------------------------
% PHẦN 4: AI PIPELINE & THUẬT TOÁN
% -------------------------------------------------------------
\section{Thuật toán Cốt lõi}

\begin{frame}{5. Trái tim Hệ thống: 5 Lớp của AI Optimizer Pipeline}
    Mục tiêu: Dập tắt \textit{"Sự bùng nổ không gian trạng thái"} (State Space Explosion).
    \vspace{0.2cm}
    \begin{enumerate}
        \item \textbf{Graph Ingestion Layer:} Dựng DAG, tính Forward/Backward Pass (Early Start, Late Finish, Critical Path).
        \item \textbf{Pruning Layer (Cắt tỉa):} Lõi tối ưu giảm tải phép tính.
        \item \textbf{Action Generation Layer:} Sinh kịch bản Đơn \& Kép.
        \item \textbf{Validation Layer (Beam Search):} Giữ lại Top K kịch bản khả thi.
        \item \textbf{Evaluation Layer:} Tính điểm Sinh lời (ROI).
    \end{enumerate}
\end{frame}

\begin{frame}{6. Pruning Layer (Lớp Cắt tỉa không gian tìm kiếm)}
    Đây là kỹ thuật mấu chốt giúp AI chạy mượt mà trên các dự án khổng lồ:
    \vspace{0.3cm}
    \begin{itemize}
        \item \textbf{Hard Pruning (Critical Path Pruning):} 
        Lập tức VỨT BỎ toàn bộ các tasks có \texttt{Total Float > 0} khỏi danh sách ứng viên rút ngắn tiến độ. Ép tiến độ task rảnh rỗi không mang lại lợi ích gì.
        \item \textbf{Soft Pruning (Hàm Phạt - Penalty Function):} 
        Chấp nhận dời lịch task trên đường găng (\texttt{Float = 0}) nếu Tiền phạt trễ dự án (Delay Penalty) \textit{rẻ hơn} chi phí thuê ngoài giờ (Crashing Cost) để bảo vệ Dòng tiền (Cashflow).
    \end{itemize}
\end{frame}

\begin{frame}{7. Nghệ thuật Đánh đổi (Constraint Relaxation)}
    Tầng 3 của dữ liệu cho phép AI \textbf{"Mở khóa/Nới lỏng ràng buộc"} khi kịch đường cùng.
    \vspace{0.2cm}
    \begin{block}{Kịch bản: Kích hoạt lệnh AGENDA\_OVERRIDE}
        \begin{itemize}
            \item \textbf{Hành động:} Tự động sửa biến \texttt{weekly\_schedule}, biến Thứ 7 \& CN thành Working Days.
            \item \textbf{Chi phí:} Hệ số làm ngoài giờ (\texttt{overtime\_multiplier} x1.5) kích hoạt, đẩy \texttt{Resource Cost} lên cao. \texttt{Fixed Cost} bất biến.
            \item \textbf{Lợi ích:} Kéo giãn tiến độ thực tế (Calendar Days) mà không làm "vỡ" giới hạn \texttt{Max Availability} của ngày thường.
        \end{itemize}
    \end{block}
\end{frame}

% -------------------------------------------------------------
% PHẦN 5: OUTPUT
% -------------------------------------------------------------
\section{Kết quả}

\begin{frame}[fragile]{8. Output Cuối cùng: AI Insight JSON Schema}
    Toàn bộ giải pháp tối ưu được gói gọn vào API Contract trả về cho Frontend:
    \vspace{0.2cm}
\begin{lstlisting}[language=json]
{
  "action_type": ["CRASHING", "AGENDA_OVERRIDE"],
  "target_tasks": ["T12", "T15"],
  "human_message": "Mo khoa lich cuoi tuan cho Cum móng. Rut ngan 4 ngay. Chi phi tang $2000.",
  "impact": {
    "time_saved_days": 4,
    "cost_penalty_usd": 2000,
    "roi_score": 0.85
  },
  "risk": {
    "float_consumed_days": 2,
    "peak_resource_variance": 0.15
  }
}
\end{lstlisting}
\end{frame}

% Slide Kết
\begin{frame}
    \centering
    \Huge \textbf{CẢM ƠN HỘI ĐỒNG ĐÃ LẮNG NGHE!}
    
    \vspace{1cm}
    \Large \textit{Q\&A}
\end{frame}

\end{document}
